\section{Conclusion}\label{sec-10}

We introduced a multi-pillar agent-native benchmark for rare disease
diagnosis spanning five capability pillars, three diagnostic data layers,
eight published systems (including two classical/offline baselines) plus one
no-scaffold control, and four hosted backbones where compatible. H1--H11 and A1--A12
are repository-defined rather than formally pre-registered. The benchmark surfaces five
findings with concrete reviewer-defensible numbers. First,
classical/offline approaches (LIRICAL Bayesian, VC-RDAgent IC+Poincar\'e)
decisively lead on Phenopacket-Store (0.47 vs best-LLM 0.30, a 17 pp
gap) and at the super-rare prevalence tail, while RareBench is near parity
(0.30 vs 0.28). They consume no hosted-LLM tokens. This is the most consequential single
result, and motivates our classical baseline column as a permanent
part of any future rare-disease agent leaderboard. Second,
multi-agent scaffolding helps only marginally (\ensuremath{\approx}0--2 pp R@1, within
overlapping CIs) in its best cells and is often neutral or harmful when the
scaffold's design (OSCE dialogue, panel orchestration) is mismatched to the
input modality. Third, DeepSeek V4-Flash is approximately 10\ensuremath{\times} cheaper
than Gemini Flash in the frozen diagnostic receipts (\$0.33 vs \$3.44 per
1,000 successful predictions) but usually trades off accuracy (\ensuremath{-}3 to \ensuremath{-}8 pp on Phenopacket-Store
and \ensuremath{-}4 to \ensuremath{-}14 pp on RareArena; RareBench is mixed) ---
the \emph{cost-efficient}, not quality-equivalent, choice for rare-disease
deployment. Fourth, GPT-5 with \texttt{reasoning\_effort=minimal} is the
costliest backbone with no consistent accuracy edge (near-competitive on
MedAgents, about \ensuremath{-}9 pp on AgentClinic dialogue), demonstrating that frontier reasoning
models are brittle when their core mechanism is disabled. Fifth, whether
faithfulness ranks decouple from accuracy ranks is judge-dependent and we report it as exploratory (same-trace Spearman \ensuremath{\rho} = 0.457 under a Gemini judge,
0.640 under a Claude judge; the repository-plan \ensuremath{\rho} \textless{} 0.5 threshold is met under one
judge but not the other) --- the durable, direction-independent point is that
accuracy-only evaluation can undersell the risk profile of rare-disease AI, since
a correct diagnosis can still rest on an unfaithful reasoning trace.

We frame these findings as retrospective decision support evaluation, not autonomous diagnosis. No clinical deployment claims are made; the
benchmark, harness, adapter shims, per-cell receipts, and leaderboard
are released to enable independent replication and to ratchet up the
shared evaluation standard in this rapidly-growing area.

\textbf{Limitations} are detailed in \Cref{sec-9}; the principal three are deferred
Pillar 4 (family-aware) reporting, single-language English evaluation,
and the LLM-generated silver-gold reference for Pillar 1 (mitigated by
physician annotation planned for the 198-case post-cutoff holdout).

\textbf{Reproducibility}: the frozen diagnostic cost receipt aggregates 90,046
successful attempts; released per-cell artifacts include OpenRouter
request IDs where available, adapter invocation metadata, and environment
records. The repository analysis plan and its unregistered OSF draft are
described transparently in \Cref{sec-5-4}; no formal pre-registration claim is made.

